\section{Introduction}

% P1: The known
Network-based drug prioritization maps compound targets onto the human interactome and scores their proximity to disease-associated protein modules \citep{Hopkins2008, Barabasi2011, Menche2015, Guney2016}. The field standard, the proximity Z-score, calibrates each compound against a size- and degree-matched null and answers the question: is this compound unusually close to this disease? Guney et al.\ \citep{Guney2016} validated this Z-score as a per-pair classifier across 402 known drug--disease associations (mean 3.5 targets per drug) and showed that, within each drug's own null, the standardization removes the dependence that raw distance has on target count and degree.

% P2: The gap
Comparing Z-score \textit{magnitudes} across compounds, however, introduces a distinct requirement: the Z-score must be monotone in the underlying topological effect. When target counts differ sharply, the larger set's null standard deviation shrinks as $|T|^{-1/2}$---a consequence of the law of large numbers on an averaged statistic---so a compound with weaker raw proximity can register stronger standardized evidence. The American Statistical Association emphasizes that statistical significance ``does not measure the size of an effect or the importance of a result'' \citep{Wasserstein2016}, and whether this non-monotonicity materializes in network-medicine practice, and under what conditions, has not been systematically examined.

% P3: Our approach
We address this with a controlled stress test in a liver-expressed interactome. We select two \textit{Hypericum perforatum} constituents---Hyperforin (10 targets) and Quercetin (62 targets)---as a deliberately high-contrast diagnostic pair. Hyperforin is the PXR-activating, CYP-inducing constituent responsible for St John's Wort's clinically important hepatic drug--drug interactions \citep{Moore2000, Watkins2001, Chen2022}, making it a mechanistically motivated comparator for engagement of the DILI module's xenobiotic-metabolism axis; Quercetin provides the high-target-count comparator. The drug-induced liver injury (DILI) module serves as the disease reference, chosen because its inclusion of cytochrome-P450, transporter, and nuclear-receptor genes ---genes central to xenobiotic metabolism, transport, and regulatory response---exposes a second question: when does direct target--disease overlap inflate apparent network influence? We complement the pair with a degree-controlled calibration benchmark ($20{,}000$ probes per size, $500{,}000$ cross-size pairs) that characterizes the operating regime in which rank reversal can occur, and we introduce perturbation efficiency---mean per-target random-walk influence on the disease module---as a complementary effect-size measure. Target provenance, botanical context, and chemical-similarity controls are detailed in Methods and \S\ref{subsec:chemsim}.

% P4: What we found (roadmap)
We show that proximity Z-score magnitude and raw closest distance can dissociate under target-count asymmetry, trace this to the $|T|^{-1/2}$ null-precision law, and confirm that the law holds across shortest-path, random-walk, and expression-weighted influence alike. A degree-controlled liver-network benchmark reveals that material rank reversal is a located, conditional outcome---rare for generic target sets, realized when the smaller set is above the 90th percentile of probe-pair margins. On the diagnostic pair, perturbation efficiency ranks the PXR/CYP-inducing constituent higher while the direct--propagated decomposition exposes the direct target--DILI overlap that an uncritical headline number would have conflated with propagated influence. This is a controlled audit of how size-dependent null precision affects the interpretation of network-proximity rankings, not a DILI-prediction model.
